% Mathématiques 4e — cours V3
% Calcul littéral et équations

\section{Objectifs}

\begin{itemize}
\item Comprendre la structure d'une expression littérale.
\item Substituer une valeur à une lettre.
\item Reconnaître une identité.
\item Développer un produit par distributivité simple.
\item Factoriser une somme.
\item Réduire une expression.
\item Tester si un nombre est solution d'une équation.
\item Formaliser les notions d'inconnue et de solution.
\item Mettre un problème en équation.
\item Résoudre algébriquement une équation du premier degré simple.
\end{itemize}

\section{1. Expression littérale}

Une expression littérale contient une ou plusieurs lettres représentant des nombres.

Exemples :
\[
3x+5,\qquad 2(a-4),\qquad L\ell.
\]

La lettre peut représenter une variable, une grandeur ou une inconnue selon le contexte.

\section{2. Structure d'une expression}

L'opération principale détermine la structure.

\[
3x+5
\]
est une somme.

\[
4(x+2)
\]
est un produit.

Cette distinction est essentielle pour choisir une transformation adaptée.

\section{3. Substitution}

Pour :
\[
A=3x^2-2x+4
\]
avec :
\[
x=2,
\]
on écrit :
\[
A=3\times2^2-2\times2+4.
\]

Puis :
\[
A=12-4+4=12.
\]

\section{4. Identité}

\begin{definition}
Une identité est une égalité vraie pour toutes les valeurs des lettres pour lesquelles les expressions ont un sens.
\end{definition}

Par exemple :
\[
3(x+2)=3x+6
\]
est une identité.

Au contraire :
\[
3(x+2)=3x+2
\]
est fausse en général.

\section{5. Distributivité}

\coursefigure{/images/mathematiques/4eme/litteral-distributivite.svg}{Rectangle découpé illustrant l'identité $k(a+b)=ka+kb$.}{Développer et factoriser sont deux lectures de la même propriété de distributivité.}

\begin{propriete}
\[
\boxed{k(a+b)=ka+kb}
\]
et :
\[
\boxed{k(a-b)=ka-kb}.
\]
\end{propriete}

\section{6. Développer}

\[
5(2x+3)=10x+15.
\]

De même :
\[
-3(x-4)=-3x+12.
\]

Le facteur extérieur doit multiplier chacun des termes de la parenthèse.

\section{7. Factoriser}

\[
8x+24=8(x+3).
\]

On met en évidence le facteur commun :
\[
8.
\]

\section{8. Réduire}

\[
3x+7x-4=10x-4.
\]

On regroupe les termes de même nature.

En revanche :
\[
10x-4
\]
ne se réduit pas davantage.

\section{9. Développer puis réduire}

\[
B=4(2x-3)+5x+1.
\]

On développe :
\[
B=8x-12+5x+1.
\]

Puis on réduit :
\[
\boxed{B=13x-11}.
\]

\section{10. Tester une égalité}

Pour :
\[
2x+5=17,
\]
on teste :
\[
x=6.
\]

Le membre de gauche vaut :
\[
2\times6+5=17.
\]

Donc $6$ est solution.

Avec :
\[
x=5,
\]
on obtient :
\[
15\neq17.
\]

\section{11. Équation}

\begin{definition}
Une équation est une égalité contenant une inconnue.

Résoudre une équation consiste à trouver toutes les valeurs de l'inconnue qui rendent l'égalité vraie.
\end{definition}

\section{12. Principe d'équilibre}

\coursefigure{/images/mathematiques/4eme/equation-balance.svg}{Balance représentant l'égalité $3x+5=26$.}{Une même transformation effectuée sur les deux membres conserve l'égalité.}

Si :
\[
A=B,
\]
ajouter le même nombre aux deux membres conserve l'égalité.

On peut aussi soustraire le même nombre, ou multiplier/diviser les deux membres par un même nombre non nul.

\section{13. Équation x+b=c}

\[
x+7=19.
\]

On soustrait $7$ :
\[
x=12.
\]

Vérification :
\[
12+7=19.
\]

\section{14. Équation ax=c}

\[
5x=35.
\]

On divise les deux membres par $5$ :
\[
\boxed{x=7}.
\]

\section{15. Équation ax+b=c}

\[
3x+5=26.
\]

On soustrait $5$ :
\[
3x=21.
\]

Puis on divise par $3$ :
\[
\boxed{x=7}.
\]

\section{16. Inconnue dans les deux membres}

\[
5x+2=3x+14.
\]

On soustrait $3x$ :
\[
2x+2=14.
\]

Puis :
\[
2x=12.
\]

Donc :
\[
\boxed{x=6}.
\]

Vérification :
\[
5\times6+2=32
\]
et :
\[
3\times6+14=32.
\]

\section{17. Coefficients décimaux}

\[
2{,}5x=15.
\]

On divise par :
\[
2{,}5.
\]

Donc :
\[
\boxed{x=6}.
\]

\section{18. Mise en équation}

Un abonnement comprend $8\ \text{€}$ de frais fixes et $3\ \text{€}$ par séance.

La facture vaut :
\[
35\ \text{€}.
\]

Si $x$ est le nombre de séances :
\[
3x+8=35.
\]

Alors :
\[
3x=27
\]
puis :
\[
x=9.
\]

\section{19. Programme de calcul}

Programme :
\begin{enumerate}
\item choisir un nombre ;
\item le multiplier par $4$ ;
\item soustraire $7$.
\end{enumerate}

La sortie est :
\[
4x-7.
\]

Si la sortie vaut $29$ :
\[
4x-7=29.
\]

Donc :
\[
x=9.
\]

\section{20. Démontrer}

Montrons que la somme de trois entiers consécutifs est un multiple de $3$.

On note :
\[
n,\quad n+1,\quad n+2.
\]

La somme vaut :
\[
3n+3=3(n+1).
\]

Elle est donc multiple de $3$.

\section{21. Contre-exemple}

L'affirmation :
\[
x^2>x\text{ pour tout }x>0
\]
est fausse.

Prenons :
\[
x=\dfrac12.
\]

Alors :
\[
x^2=\dfrac14<\dfrac12.
\]

Un seul contre-exemple suffit.

\section{22. Tableur et conjecture}

Un tableur permet de tester rapidement une expression sur de nombreuses valeurs.

Cela aide à formuler une conjecture.

Mais une collection d'exemples ne remplace pas une démonstration générale.

\section{23. Géométrie et équation}

Un rectangle a pour longueur :
\[
x+3
\]
et largeur :
\[
x.
\]

Son périmètre vaut $30$ :
\[
2(x+3)+2x=30.
\]

On développe :
\[
4x+6=30.
\]

Donc :
\[
x=6.
\]

\section{24. Méthode de résolution}

\begin{methode}
\begin{enumerate}
\item choisir l'inconnue ;
\item écrire l'équation ;
\item simplifier chaque membre ;
\item appliquer les mêmes transformations aux deux membres ;
\item isoler l'inconnue ;
\item vérifier ;
\item interpréter la solution.
\end{enumerate}
\end{methode}

\section{25. Résoudre un problème avec deux étapes de modélisation}

Un cinéma propose une carte coûtant :
\[
18\ \text{€}
\]
puis chaque place coûte :
\[
6\ \text{€}.
\]

Sans carte, chaque place coûte :
\[
9\ \text{€}.
\]

On cherche à partir de combien de séances la carte devient avantageuse.

Avec la carte :
\[
C=18+6x.
\]

Sans carte :
\[
S=9x.
\]

On cherche :
\[
18+6x<9x.
\]

Donc :
\[
18<3x,
\]
puis :
\[
x>6.
\]

À partir de :
\[
\boxed{7\text{ séances}},
\]
la carte devient moins chère.

Cet exemple montre qu'une équation ou une inéquation simple est d'abord la traduction d'une situation.

\section{26. Tester une formule générale}

Considérons l'expression :
\[
A=(x+3)^2-x^2.
\]

On peut développer :
\[
(x+3)^2=(x+3)(x+3).
\]

Donc :
\[
A=x^2+6x+9-x^2.
\]

Ainsi :
\[
\boxed{A=6x+9}.
\]

Tester avec :
\[
x=2
\]
donne :
\[
(2+3)^2-2^2=25-4=21,
\]
et :
\[
6\times2+9=21.
\]

L'essai confirme le calcul pour cette valeur, mais c'est le développement qui démontre l'identité pour toutes les valeurs de $x$.

\section{27. Organiser une démonstration littérale}

Pour démontrer une propriété avec une lettre :
\begin{methode}
\begin{enumerate}
\item choisir une lettre pour représenter un nombre quelconque ;
\item traduire les autres nombres ou grandeurs avec cette lettre ;
\item écrire l'expression correspondant à la propriété ;
\item développer, réduire ou factoriser ;
\item faire apparaître une forme caractéristique ;
\item conclure en revenant à la propriété à démontrer.
\end{enumerate}
\end{methode}

Par exemple, un entier pair s'écrit :
\[
2n.
\]

La somme de deux entiers pairs :
\[
2n+2p
\]
se factorise :
\[
2(n+p).
\]

Elle est donc encore paire.

Le calcul littéral permet ici de dépasser les exemples numériques pour obtenir une conclusion générale.

\section{28. Erreurs fréquentes}

\begin{itemize}
\item Oublier de distribuer un facteur à tous les termes.
\item Réduire des termes de nature différente.
\item Modifier un seul membre d'une équation.
\item Changer un signe sans effectuer d'opération justifiée.
\item Confondre identité et équation.
\item Croire que quelques essais démontrent une identité.
\item Oublier de vérifier ou d'interpréter la solution.
\end{itemize}

\section{À retenir}

\begin{aretenir}
La distributivité donne :
\[
\boxed{k(a+b)=ka+kb}.
\]

Développer transforme un produit en somme ; factoriser transforme une somme en produit.

Une équation est une égalité contenant une inconnue.

Pour la résoudre, on effectue des transformations qui conservent l'égalité jusqu'à isoler l'inconnue.

Le calcul littéral sert aussi à démontrer des propriétés générales.
\end{aretenir}
